\section{Discussion}
\label{sec:discussion}

\subsection{The hypotheses}
\label{sec:verdicts}

\begin{table}[tbp]
  \centering
  \small
  \begin{tabular}{@{}>{\raggedright\arraybackslash}p{0.2\linewidth}
                     >{\raggedright\arraybackslash}p{0.17\linewidth}
                     >{\raggedright\arraybackslash}p{0.55\linewidth}@{}}
    \toprule
    Hypothesis & Verdict & Evidence \\
    \midrule
    H1, planning & holds for Reynolds &
      Reynolds 0 $\to$ 8 and 0 $\to$ 10 of 10; hybrid and Couzin within 4
      seeds. The ``no effect where the way is clear'' half was not tested. \\
    \addlinespace
    H2, the corner trap & holds &
      Reynolds never got round with the plan off, in tuning or evaluation;
      with the plan, 10 of 10 on \texttt{l\_turn}. \\
    \addlinespace
    H3, tuned Couzin & half &
      Reaches C as often as the hybrid; its drones touch as the hybrid's do,
      and it is the only model that crashed. \\
    \addlinespace
    H4, leaders & Couzin only &
      Couzin 7, 16, 28 of 30 for 1, 2, 4 leaders; Reynolds flat; the hybrid
      worst with 2. \\
    \bottomrule
  \end{tabular}
  \caption{The hypotheses of \cref{sec:hypotheses} against the results.}
  \label{tab:verdicts}
\end{table}

\Cref{tab:verdicts} closes the loop. Two hypotheses hold, one half holds,
and the research question of the project, H4, holds for one model out of
three. The rest of this section argues that these are not three separate
outcomes but one.

\subsection{What drives a follower}
\label{sec:drive}

The three models differ in many terms, but one difference explains most of
the results: what makes a follower move.

\begin{description}
  \item[Hybrid: its own cruise.] A follower cruises at full speed along its
    neighbours' mean heading. The flock therefore has momentum. On
    \texttt{l\_turn} that momentum carries the leaders round the corner even
    with no plan, where Reynolds' leaders stop (10 of 10 against none). On
    \texttt{dead\_end} the same momentum is the problem: two leaders heading
    round opposite arms of the cup pull the mean heading two ways, the pulls
    cancel, and the flock cruises on between them, into the cup. The flock is
    quick to follow a heading, and so it follows the average one.
  \item[Reynolds: nothing but its neighbours.] A follower has no drive of its
    own; it moves only as its neighbours pull it. With leaders that steer
    straight at C, that is useless: the leaders stop at a wall and the flock
    stops with them. With leaders that follow a plan, it is faithful: the
    flock speeds up slowly, the leaders get ahead before it commits, and when
    two leaders disagree the flock goes with one of them. Reynolds' success
    hardly depends on the number of leaders, presumably because one leader
    pulling is enough for a flock that has nothing pulling it the other way.
  \item[Couzin: a direction to be outvoted.] A follower keeps a direction as
    its own state and turns it toward its neighbours'. Leaders have to outvote
    the flock, so more leaders work better, exactly as
    \citet{couzin2005} predicted; one leader among 30 is lost. And since a
    Couzin leader feels its flock only while it still sees it, a leader that
    gets ahead can leave the flock behind for good.
\end{description}

So no model is simply best. Each one's strength and its weakness come from the
same property. Momentum makes the hybrid robust to leaders that steer badly
and fragile to leaders that disagree; having no drive makes Reynolds' flock
useless without a plan and faithful with one; state makes Couzin's flock
need enough leaders and lets its leaders run away. Which model to choose
depends on what the leaders can be trusted to do.

\subsection{A plan helps only if the flock follows}
\label{sec:plan-and-flock}

The planning layer was designed for the leaders alone, and on that level it
works: in the replays, no leader stayed stuck in the cup once it could plan
again. But a better path for the leaders helps only as far as the flock
follows them, and the model decides that. The leaders also plan
independently, as if they were interchangeable: each from where it stands,
each for itself. This was expected to split the flock at worst, each part with
a leader. For a model that averages its leaders, two independent plans are
worse than one: the hybrid succeeds in 7 of 10 runs whose leaders happen to
choose the same side, and in 2 of 20 whose leaders don't.

The deliberative layer is therefore missing a part. A hybrid architecture
plans for the system it controls, and here the system is the flock, not
the leader. Leaders that agree on one path, or that plan with the flock's
response in mind, are the natural next step (\cref{sec:conclusion}). Without
being recognised, they could still agree: they share the map and the goal,
so a rule as simple as every leader planning from the same point, such as
A's centre if the map gave it to them, would give one path.

\subsection{Reliability against safety}
\label{sec:safety}

Reynolds with the plan is the most reliable model on \texttt{dead\_end} with 20
drones and 2 leaders, and it keeps its drones about 0.5\,m apart, never closer than
0.33\,m. The hybrid and Couzin reach touching distance, 0.10\,m, in many runs,
and Couzin crashed in 4 runs of 40. In simulation that costs little; on real
drones it would matter more. Webots does not simulate propeller downwash, which
makes a drone flying over another unstable, and real IR distances are noisier
than Webots' (1.4\,cm here). The stopping limit kept every drone off the
walls but did not keep them apart in a crowd: it is one-sided, each drone
limits only itself, and the flight layer follows it a time constant late.
\citet{soria2021} reach a similar conclusion about reactive swarm models in
clutter: they don't guarantee safe flight, which is why they replaced them
with model predictive control. For a real flock the spacing would weigh as
much as the success rate.

\subsection{Back to the literature}
\label{sec:literature}

\begin{itemize}
  \item \citet{couzin2005} found that a group with two informed subgroups
        follows their average when the two preferred directions differ
        little, and picks one when they differ a lot. The hybrid shows the
        first behaviour on drones, at the start of two paths that part slowly;
        Reynolds' flock, which commits later, shows the second.
  \item Larger flocks needed more leaders in every model. That is consistent
        with Couzin's result, which is about the \emph{fraction} of leaders:
        it falls for larger groups, but not as fast as $1/n$, so the number
        of leaders needed still grows.
  \item \citet{vasarhelyi2018}'s model was built for flocks that stay inside
        an arena with no goal: there, walls that push the flock back are the
        point. For a flock that has to get somewhere, through a forest and a
        gap, the push became a bounce, and the stopping limit replaced it.
  \item \citet{viragh2014} described the robot by its imperfections; the
        first flight of Reynolds' rules failed exactly as they warn, and the
        fix was to make the robot follow, not to change the rules.
\end{itemize}

\subsection{Limitations}
\label{sec:limitations}

\begin{itemize}
  \item \emph{Few seeds, many tests.} Ten seeds per setting resolve large
        effects only, and the $p$ values are uncorrected. Effects of 2--4
        seeds, such as the plan for the hybrid and Couzin, are not resolved
        either way.
  \item \emph{One trap map.} The second experiment ran on \texttt{dead\_end}
        only, and \texttt{l\_turn} sits at the ceiling. The leader-split
        mechanism depends on a symmetric obstacle with two ways round; another
        trap might rank the models differently.
  \item \emph{Unequal tuning.} The hybrid was tuned over a whole development
        phase, Couzin and Reynolds with 15 runs each. The held-out maps limit
        the damage but don't remove it.
  \item \emph{Simulation.} Flat flight only, idealised sensing (IR without
        noise or occlusion beyond Webots' own, obstacles known from the map),
        and no downwash.
  \item \emph{The split analysis} infers each leader's way round from a plan
        made at its spawn, a stand-in for the plan it made at go time; the
        replays agree with it, but it was not logged for every run. Why an
        even 2--2 split with 4 leaders does not break the hybrid is not
        explained.
\end{itemize}
